% !TEX root = ../nc-cubic.tex
\section{残余圏の形式認識と到達範囲}
\label{real:sec:scope}
\subsection{通常 Artin 変形の認識}
\begin{maintheorem}[残余圏の通常の Artin 形式認識]
滑らかな三次 $X$ の dg 残余圏を $\mathcal K_X=\operatorname{Ku}(X)$ とし、
$\mathcal D_X=\operatorname{Ind}(\mathcal K_X)$ とする。
$n\ge3$、$n\ne4$ なら、特殊ファイバーを同定した通常の可換局所 Artin 基底上で
\[
 \operatorname{Def}_X\simeq
 \tau_{\le1}\operatorname{CatDef}^{c}_{\mathcal D_X}.
\]
両辺の同型類は $\binom{n+2}{3}$ 次元の滑らかな形式円板となる。
$n=4$ では圏変形の形式 hull の通常 Artin 同型類関手が
二十二次元形式円板となり、幾何的変形の円板は二十次元である。
\end{maintheorem}

証明は定理\ref{thm:nf-residual-formal}に与える。
既知の sector 公式と幾何的 Kodaira--Spencer 単射性から得る接写像同型を、
実際の幾何的変形の全次数の自然変換と合わせる。
$HH^3$ の消滅に依存しないため、$HH^3=\C^2$ となる $n=7$ も含む。
圏変形の形式 hull と実際の compactly generated 変形の比較には
Lurie \cite{NFDAGX} の tame 圏に対する結果を用いる。
未切断の高次群oidや導来 Artin 基底上での同値は主張しない。

rotation $\mathsf R$ と Serre 関手 $S$ の既知の関係は
$\mathsf R^3=[2]$、$S=\mathsf R^{-(n-1)}[n]$ である。
$n\not\equiv1\pmod3$ では rotation を Serre と shift から回復できるが、
$n\equiv1\pmod3$ ではこの方法は使えない。
これは四次元の形式認識の例外とは異なる区分である。


\subsection{相対比較と降下の位置付け}

構造対象付きモジュライ、trace--Hessian と残余三次の差、
点曲線による尺度の規格化、比較核の Picard 障害と周期位相は、
モジュライ・降下稿\cite{NCCubicModuliDescent}で扱う。
自然な標的を指定する比較や族上の比較では、本稿からもそこでの
中心直線・相対尺度・scalar 自明化の結果を参照する。
固定された複素係数で任意の平面方程式を選ぶ直接 Ku 比較は、
これらの精密な降下条件を前提としない。
幾何学的な同値の存在、核の同型類の定義体、実際の核の scalar 障害、
rotation と指定周期の位相は、それぞれ異なる問題である。

\subsection{可換実現の境界と残る問題}

中心課題は主予想\ref{conj:nc-main}であり、
既知の可換実現 locus の境界延長と各 intrinsic 成分を覆う表示を優先する。
点スキームと Poisson 零スキームの一致を一般に期待できないことは、
注意\ref{ppz:rem:naive-expectation}とそこで参照する三様体・五様体の例で説明した。
既知の分離型における一致とその限界は、可換実現の幾何を理解するために保持する。
中心変数を持つ分離型の可換実現は定理\ref{csc:thm:comparison}で扱い、
定数型混合族では、有限中心性と相対対角の証明から
系\paperref{mixed}{cmo:cor:cofinite}の余有限可換実現を得た。
指定二因子では、定理\paperref{mixed}{cmo:thm:morse-singular-sixty-four}により
norm と交わらない六十四係数で全 coherent Ku の非 proper 性が確定した。
定理\paperref{mixed}{cmo:thm:fourteen-node-orders}で十四節点の局所代数と全単純加群を決定し、
中央切断・相対対角・射影境界の証明を合わせて
定理\paperref{mixed}{cmo:thm:exact-sixty-four}の正確な非 smooth 集合 $V(P_{64})$ を得た。
例\paperref{mixed}{cmo:ex:sixty-four}の norm の全六十四係数と $c=0$ は smooth proper である。
定理\paperref{mixed}{cmo:thm:all-central-commutative-kernel}は、
この指定族の全 smooth locus 上に、明示可換三次族への相対 perfect 核を構成する。
norm の全六十四根にも同じ核が適用される。
具体的な比較先として、定理\paperref{mixed}{cmo:thm:commutative-discriminant-candidate}の
可換三次族と、零係数での命題\paperref{mixed}{cmo:prop:zero-eckardt-resolution}の模型を使う。
零係数では定理\paperref{mixed}{cmo:thm:eckardt-root-residual}の根スタック上の十六直交条件と
指定された指標作用を比較核の標的として使える。
さらに定理\paperref{mixed}{cmo:thm:root-order-residual}の非可換側の実二成分分解と、
境界上の twist 射および高次構造の比較から、
定理\paperref{mixed}{cmo:thm:fixed-zero-commutative-kernel}の
$\Ku_{\mathrm{coh}}(A_0)\simeq\Ku(V(G-u^2L))$ の実 perfect 核を得る。
その構成は $\C$ 上の通常の dg Morita 同値であり、
全標識の同時保存や有理数体への核降下までは主張しない。
全中心係数の構成では、次数一の中心超平面制限を使う。
その comonad の二項・三項の構造を混合係数と中心係数へ同定し、
四項以降の比較の一意性から実際の Morita 同値族を得る。
特異係数を含む残余線束圏の比較が先にあり、
既知の判別集合はその後の尺度決定に用いる。
系\paperref{mixed}{cmo:cor:all-alpha-central-kernel}は
固定した二つの Sklyanin 因子の任意の定数補正も扱う。
交換子パラメータ $h$ については、定理\paperref{mixed}{mix:thm:h-exact-candidate}により
指定した可換候補への全形式次数の MF 同値と一般点の実 perfect 核を得た。
$h=0,1$ を含む相対 smooth proper 模型と、
$h=1$ の DVR 上の perfect 整延長は
定理\paperref{mixed}{mix:thm:h-semilocal-smooth}と
\paperref{mixed}{mix:thm:h-integral-kernel}で構成できる。
選択済み $h$ 一般点核を特殊点でも可逆に延長し、
固定 $h=1$ の核と相対的に整合させることは残る。
通常の Hochschild 群と有理 Euler 形式は全 smooth 係数で一致するが、
それだけで比較核は得られない。
非次数付き affine MF 比較は同領域で零となるため、
本稿の実比較は、内部次数と実随伴を持つ射影残余圏から構成している。
他の族や選択済み一般点核に対しては、
命題\ref{grl:prop:trait-kernel-test}の単位・余単位の有限長判定が引き続き必要となる。
通常の周囲代数の derived Morita 比較と、形式的な $h$ 進核の
$h=1$ への無条件な代入は、この手順の代わりにならない。
可換実現 locus の開性だけでは全 smooth locus への拡張は従わない。
次の中心課題は、他の Sklyanin パラメータや線形補正、
複数の中心変数を含む混合表示への実比較の拡張と、全 intrinsic 成分の被覆である。

精密な降下の側には、周期位相の他の境界での分岐、正準位相体上の scalar index、
代数的核の coherent な中心線型性、高次元の任意の full class と非分裂軌道が残る。
これらの記法と条件はモジュライ・降下稿に置く。
特に同稿の補題\paperref{moduli}{lem:nfm-period-two-boundary-criterion}は、
一つの完備 trait 上の位相値だけから大域式を結論できないことを明記する。
同稿の第\paperref{moduli}{subsubsec:nfm-boundary-products-period}項では、
全 orbit Hom と指定周期を保つ整模型で低次積・飽和関係・三次積の消失次数を調べる。
この積の計算と局所尺度の比較は、滑らかな残余模型や full 比較核自体の
境界延長とは別である。
形式座標の代数性と内在的なモジュライの代数性も区別し、
境界や全 Abel 形式近傍は、一般次元の可換実現を妨げる障害の解消に必要な範囲で扱う。
